% Section 9.2, from lectures/L09/README.md: the goals, the evaluation questions, and the next
% lecture.

\needspace{10\baselineskip}
\section{Review}
\label{c9:sec:review}

\subsection*{What you should be able to do now}
After this chapter, you should:
\begin{itemize}
  \item Explain why hardcoded per-instance functions or macros break down once a program needs
    more than one of something.
  \item Be able to design a struct that bundles a peripheral's configuration and state.
  \item Be able to write ``methods'' as plain functions taking a pointer to that struct as their
    first parameter (the \code{self} convention).
  \item Refactor a program built on duplicated globals into one built on an array of struct
    instances.
\end{itemize}

\needspace{6\baselineskip}
\subsection*{Questions to test yourself}
\begin{enumerate}
  \item What problem does bundling a pin number into a struct solve that a \code{#define}d macro
    can't?
  \item Why does \code{gpio_read()} take a pointer to a \code{gpio_t} rather than a \code{gpio_t}
    by value?
  \item How would you rewrite five separate
    \code{LED1_ON}/\allowbreak\code{LED2_ON}/\allowbreak\dots{} macros as a single
    \code{gpio_write()} function plus an array of \code{gpio_t} instances?
\end{enumerate}

\needspace{6\baselineskip}
\subsection*{Looking ahead}
Chapter~\ref{c10:ch}: hiding a struct's implementation details (opaque structs) and manual
polymorphism via function pointers and vtables.
